% ============================================================
% PROPOSAL PROYEK AKHIR — MUHAMMAD DIAS AL-IZZAT
% PENGEMBANGAN SISTEM KLASIFIKASI SKETSA BERBASIS CNN MOBILENET
% DENGAN CONFIDENCE SCORE DAN CLUSTERING POLA KEPUTUSAN
% PADA SISTEM LITERASI KECERDASAN BUATAN UNTUK SISWA SMP
%
% Format: PENS Standard (A4, 12pt TNR, 1.5 spacing, 4cm/3cm margin)
% Mengikuti pola format Izzah (PilAItes 2025)
% Skeleton file — empty sections siap diisi Phase 2
% ============================================================

\documentclass[12pt,a4paper,oneside,openany]{report}

% ============================================================
% PACKAGES
% ============================================================
\usepackage[top=4cm,left=4cm,right=3cm,bottom=3cm]{geometry}
\usepackage{graphicx}
\usepackage{booktabs}
\usepackage{array}
\usepackage{tabularx}
\usepackage{longtable}
\usepackage{float}
\usepackage{caption}
\usepackage{titlesec}
\usepackage{indentfirst}
\usepackage{setspace}
\usepackage{etoolbox}
\usepackage{fancyhdr}
\usepackage{afterpage}
\usepackage{amsmath}
\usepackage{amssymb}
\usepackage{url}
\usepackage{hyperref}
\hypersetup{colorlinks=false, pdfborder={0 0 0}}

% Times New Roman (pdflatex standard)
\usepackage{times}
\renewcommand{\rmdefault}{ptm}

% ============================================================
% LINE SPACING 1.5
% ============================================================
\setstretch{1.5}

% ============================================================
% PARAGRAPH INDENT (1.25 cm, PENS standard)
% ============================================================
\setlength{\parindent}{1.25cm}
\setlength{\parskip}{0pt}

% ============================================================
% CHAPTER / SECTION / SUBSECTION FORMATTING
% ============================================================
\titleformat{\chapter}[display]
  {\normalfont\fontsize{14pt}{16.8pt}\selectfont\bfseries\centering}
  {BAB \thechapter}
  {12pt}
  {\MakeUppercase}

\titleformat{\section}
  {\normalfont\fontsize{12pt}{14.4pt}\selectfont\bfseries}
  {\thesection}
  {1em}
  {\MakeUppercase}

\titleformat{\subsection}
  {\normalfont\fontsize{12pt}{14.4pt}\selectfont\bfseries}
  {\thesubsection}
  {1em}
  {}

\titleformat{\subsubsection}
  {\normalfont\fontsize{12pt}{14.4pt}\selectfont\bfseries\itshape}
  {\thesubsubsection}
  {1em}
  {}

\titlespacing*{\chapter}{0pt}{0pt}{24pt}
\titlespacing*{\section}{0pt}{18pt}{12pt}
\titlespacing*{\subsection}{0pt}{14pt}{8pt}

% ============================================================
% CAPTION FORMATTING (Pola Izzah)
% - Tabel: caption DI ATAS, font 11pt, label bold
% - Gambar: caption DI BAWAH, font 11pt, label bold
% ============================================================
\captionsetup[table]{font=small,labelfont=bf,name={Tabel},labelsep=period,justification=centering,aboveskip=6pt,belowskip=0pt,position=top}
\captionsetup[figure]{font=small,labelfont=bf,name={Gambar},labelsep=period,justification=centering,aboveskip=0pt,belowskip=6pt,position=bottom}

% ============================================================
% HEADER / FOOTER (page number top-right)
% ============================================================
\pagestyle{fancy}
\fancyhf{}
\fancyhead[R]{\thepage}
\renewcommand{\headrulewidth}{0pt}
\renewcommand{\footrulewidth}{0pt}

\fancypagestyle{plain}{
  \fancyhf{}
  \fancyhead[R]{\thepage}
  \renewcommand{\headrulewidth}{0pt}
}

% ============================================================
% BLANK PAGE COMMAND
% ============================================================
\newcommand{\blankpage}{%
  \afterpage{%
    \null
    \vfill
    \begin{center}
      {\small Halaman ini sengaja dikosongkan}
    \end{center}
    \vfill
    \thispagestyle{plain}
    \newpage
  }%
}

% ============================================================
% CUSTOM COMMANDS
% ============================================================
% \gbr{filename}{caption}{label}
\newcommand{\gbr}[3]{%
  \begin{figure}[H]
    \centering
    \includegraphics[width=\textwidth]{#1}
    \caption{#2}
    \label{#3}
  \end{figure}
}

% \gbrsrc{filename}{caption}{label}{source}
\newcommand{\gbrsrc}[4]{%
  \begin{figure}[H]
    \centering
    \includegraphics[width=\textwidth]{#1}
    \caption{#2}
    \label{#3}
    \begin{center}\small (Sumber: #4)\end{center}
  \end{figure}
}

% \tabel{caption}{content}{label}
\newcommand{\tabel}[3]{%
  \begin{table}[H]
    \centering
    \caption{#1}
    \label{#3}
    #2
  \end{table}
}

% ============================================================
% TOC FORMATTING
% ============================================================
\renewcommand{\contentsname}{DAFTAR ISI}
\renewcommand{\listfigurename}{DAFTAR GAMBAR}
\renewcommand{\listtablename}{DAFTAR TABEL}

% ============================================================
\begin{document}
% ============================================================

% ============================================================
% 1. COVER PAGE
% ============================================================
\begin{titlepage}
\centering
\vspace*{1cm}

{\fontsize{14pt}{16pt}\selectfont PROPOSAL PROYEK AKHIR}

\vspace{2cm}

{\fontsize{14pt}{18pt}\selectfont\bfseries\MakeUppercase{Pengembangan Sistem Klasifikasi Sketsa Berbasis CNN MobileNet dengan Confidence Score dan Clustering Pola Keputusan pada Sistem Literasi Kecerdasan Buatan untuk Siswa SMP}}

\vspace{3cm}

{\fontsize{12pt}{14pt}\selectfont Oleh:}\\[0.3cm]
{\fontsize{12pt}{14pt}\selectfont\bfseries Muhammad Dias Al-Izzat}\\[0.2cm]
{\fontsize{12pt}{14pt}\selectfont NRP. 53226000xx}

\vspace{3cm}

{\fontsize{12pt}{14pt}\selectfont Dosen Pembimbing:}\\[0.3cm]
{\fontsize{12pt}{14pt}\selectfont [Nama Dosen Pembimbing]}\\[0.2cm]
{\fontsize{12pt}{14pt}\selectfont NIP. [NIP]}

\vspace{3cm}

{\fontsize{12pt}{14pt}\selectfont PROGRAM STUDI SARJANA TERAPAN TEKNOLOGI REKAYASA MULTIMEDIA}\\[0.2cm]
{\fontsize{12pt}{14pt}\selectfont JURUSAN TEKNOLOGI MULTIMEDIA DAN JARINGAN}\\[0.2cm]
{\fontsize{12pt}{14pt}\selectfont POLITEKNIK ELEKTRONIKA NEGERI SURABAYA}\\[0.2cm]
{\fontsize{12pt}{14pt}\selectfont 2025}

\end{titlepage}

% ============================================================
% 2. BLANK PAGE
% ============================================================
\blankpage

% ============================================================
% 3. HALAMAN PENGESAHAN
% ============================================================
\chapter*{HALAMAN PENGESAHAN}
\addcontentsline{toc}{chapter}{HALAMAN PENGESAHAN}

Proposal Proyek Akhir ini disusun dan diajukan oleh:

\vspace{0.4cm}

\begin{tabular}{ll}
Nama & : Muhammad Dias Al-Izzat \\
NRP & : 53226000xx \\
Program Studi & : Sarjana Terapan Teknologi Rekayasa Multimedia \\
Jurusan & : Teknologi Multimedia dan Jaringan \\
\end{tabular}

\vspace{0.5cm}

Telah diperiksa dan disahkan oleh Dosen Pembimbing pada tanggal: \underline{\hspace{3cm}}

\vspace{1.5cm}

\begin{tabular}{ll}
Dosen Pembimbing, & \\
\\
\\
\\
\underline{\hspace{5cm}} & \\
[Nama Dosen Pembimbing] & \\
NIP. \underline{\hspace{3cm}} & \\
\end{tabular}

\blankpage

% ============================================================
% 4. ABSTRAK
% ============================================================
\chapter*{ABSTRAK}
\addcontentsline{toc}{chapter}{ABSTRAK}

% TODO Phase 2: Isi abstrak 200-250 kata + kata kunci
% Template: konteks → masalah → solusi → metode → hasil yang diharapkan
%
% [ABSTRAK — ISI DI SINI]

\vspace{0.5cm}

\noindent\textbf{Kata Kunci}: % [kata kunci 1, kata kunci 2, ...]

\blankpage

% ============================================================
% 5. ABSTRACT (English)
% ============================================================
\chapter*{ABSTRACT}
\addcontentsline{toc}{chapter}{ABSTRACT}

% TODO Phase 2: English version of abstrak
% [ABSTRACT — ISI DI SINI]

\vspace{0.5cm}

\noindent\textbf{Keywords}: % [keywords in English]

\blankpage

% ============================================================
% 6. DAFTAR ISI
% ============================================================
\chapter*{DAFTAR ISI}
\addcontentsline{toc}{chapter}{DAFTAR ISI}

\tableofcontents
\clearpage

% ============================================================
% 7. DAFTAR GAMBAR
% ============================================================
\chapter*{DAFTAR GAMBAR}
\addcontentsline{toc}{chapter}{DAFTAR GAMBAR}

\listoffigures
\clearpage

% ============================================================
% 8. DAFTAR TABEL
% ============================================================
\chapter*{DAFTAR TABEL}
\addcontentsline{toc}{chapter}{DAFTAR TABEL}

\listoftables
\clearpage

% ============================================================
% BAB 1 — PENDAHULUAN
% (UNCHANGED from existing proposal — JANGAN diubah, hanya format convert)
% Total target: ~5 halaman
% IDENTIK dengan Can — pastikan selaras
% ============================================================
\chapter{PENDAHULUAN}
\label{chap:bab1}

% ------------------------------------------------------------
\section{LATAR BELAKANG}
\label{sec:latar-belakang}
% ------------------------------------------------------------
% TODO Phase 2: Copy Bab 1.1 dari proposal-dias.tex existing
% Estimasi: ~2 halaman, 5-6 paragraf
% [ISI DI SINI]

% ------------------------------------------------------------
\section{PERMASALAHAN}
\label{sec:permasalahan}
% ------------------------------------------------------------
% TODO Phase 2: Copy Bab 1.2 dari proposal-dias.tex existing
% Estimasi: ~0.5 halaman, 1 paragraf
% [ISI DI SINI]

% ------------------------------------------------------------
\section{BATASAN MASALAH}
\label{sec:batasan-masalah}
% ------------------------------------------------------------
% TODO Phase 2: Copy Bab 1.3 dari proposal-dias.tex existing
% Estimasi: ~0.5 halaman, bullet list 4-5 poin
% [ISI DI SINI]

% ------------------------------------------------------------
\section{TUJUAN}
\label{sec:tujuan}
% ------------------------------------------------------------
% TODO Phase 2: Copy Bab 1.4 dari proposal-dias.tex existing
% Estimasi: ~0.5 halaman
% [ISI DI SINI]

% ------------------------------------------------------------
\section{MANFAAT}
\label{sec:manfaat}
% ------------------------------------------------------------
% TODO Phase 2: Copy Bab 1.5 dari proposal-dias.tex existing
% Estimasi: ~0.5 halaman, bullet 4 poin
% [ISI DI SINI]

% ------------------------------------------------------------
\section{SISTEMATIKA PENULISAN}
\label{sec:sistematika-penulisan}
% ------------------------------------------------------------
% TODO Phase 2: Copy Bab 1.6 dari proposal-dias.tex existing + REVISE
% Estimasi: ~0.5 halaman, paragraf per BAB
% REVISE: ubah mention "BAB III Perancangan Sistem" → "BAB III Deskripsi Sistem"
% + tambah mention jadwal penelitian
% [ISI DI SINI]

% ============================================================
% BAB 2 — KAJIAN PUSTAKA
% BEDA Can vs Dias — Dias fokus Backend/CNN/Data Analysis
% Total target: ~8-10 halaman
% Pola Izzah: ringkas, detail di Bab 3
% ============================================================
\chapter{KAJIAN PUSTAKA}
\label{chap:bab2}

% Paragraf pembuka bab — TODO Phase 2
% [PARAGRAF PEMBUKA BAB 2]

% ------------------------------------------------------------
\section{DESKRIPSI PERMASALAHAN}
\label{sec:deskripsi-permasalahan}
% ------------------------------------------------------------
% TODO Phase 2: Copy dari proposal-dias.tex existing + ringkas
% Estimasi: ~1 halaman, 3 paragraf
% [ISI DI SINI]

% ------------------------------------------------------------
\section{TEORI PENUNJANG}
\label{sec:teori-penunjang}
% ------------------------------------------------------------
% TODO Phase 2: Copy dari proposal-dias.tex existing + ringkas (detail di Bab 3)
% Estimasi: ~5-6 halaman total, 0.5-1 hal per teori
% Format pola Izzah: definisi → penjelasan → posisi dalam penelitian

\subsection{Literasi Kecerdasan Buatan}
\label{subsec:literasi-ai}
% [ISI DI SINI]

\subsection{Human-in-the-Loop pada Sistem Klasifikasi}
\label{subsec:hitl-klasifikasi}
% [ISI DI SINI]

\subsection{Convolutional Neural Network dan MobileNet}
\label{subsec:cnn-mobilenet}
% [ISI DI SINI]

\subsection{TensorFlow.js dan Client-Side Inference}
\label{subsec:tfjs}
% [ISI DI SINI]

\subsection{Confidence Score}
\label{subsec:confidence-score}
% [ISI DI SINI]

\subsection{Data Logging}
\label{subsec:data-logging}
% [ISI DI SINI]

\subsection{K-Means Clustering}
\label{subsec:kmeans}
% [ISI DI SINI]

% ------------------------------------------------------------
\section{PENELITIAN TERKAIT}
\label{sec:penelitian-terkait}
% ------------------------------------------------------------
% TODO Phase 2: 4 penelitian terkait + Tabel State of the Art
% Estimasi: ~2-3 halaman

\subsection{Deep Learning for Free-Hand Sketch Recognition}
\label{subsec:penelitian-1}
% [ISI DI SINI — Xu et al. [6]]

\subsection{Front-End Deep Learning Web Applications}
\label{subsec:penelitian-2}
% [ISI DI SINI — Goh et al. [7]]

\subsection{Confidence Visualization}
\label{subsec:penelitian-3}
% [ISI DI SINI — Karran et al. [9]]

\subsection{K-Means untuk Pola Belajar}
\label{subsec:penelitian-4}
% [ISI DI SINI — Pansri [10] + Alzahrani [11]]

\subsection{State of the Art}
\label{subsec:state-of-the-art}
% TODO Phase 2: Tabel 2.1 State of the Art (POLosan, format Izzah)
% Kolom: No | Penelitian | Fokus | Metode | Temuan | Celah yang Diisi
% [ISI DI SINI + Tabel 2.1]

% ============================================================
% BAB 3 — DESKRIPSI SISTEM
% BEDA Can vs Dias — DETAIL, scope Backend/CNN/Data Analysis
% Total target: ~18-22 halaman
% Urutan pola Izzah: Deskripsi Solusi → Metodologi → Desain Sistem → Jadwal
% Judul: "DESKRIPSI SISTEM" (BUKAN "PERANCANGAN SISTEM" seperti existing)
% ============================================================
\chapter{DESKRIPSI SISTEM}
\label{chap:bab3}

% Paragraf pembuka bab — TODO Phase 2
% [PARAGRAF PEMBUKA BAB 3]

% ------------------------------------------------------------
\section{DESKRIPSI SOLUSI}
\label{sec:deskripsi-solusi}
% ------------------------------------------------------------
% TODO Phase 2: Copy dari proposal-dias.tex existing Bab 3.1
% Estimasi: ~1-2 halaman
% Konten: 4 prinsip (CNN MobileNet, TF.js, REST API logging, K-Means)
% + 5 komponen teknis
% [ISI DI SINI]

% ------------------------------------------------------------
\section{METODOLOGI PENELITIAN}
\label{sec:metodologi-penelitian}
% ------------------------------------------------------------
% TODO Phase 2: PINDAH dari posisi akhir (existing 3.15) ke sini (3.2)
% Estimasi: ~2-3 halaman
% Konten WAJIB: Waterfall (SDLC) 5 tahap + Fishbone 6M + SUS + Black-Box + akurasi model
% [ISI DI SINI]

\subsection{Pendekatan dan Jenis Penelitian}
\label{subsec:pendekatan-penelitian}
% [ISI DI SINI — Mixed Methods + R&D]

\subsection{Metode Pengembangan}
\label{subsec:metode-pengembangan}
% [ISI DI SINI — SDLC Waterfall 5 tahap]

\subsection{Analisis Masalah (Fishbone 6M)}
\label{subsec:analisis-masalah}
% [ISI DI SINI — Fishbone 6M: Man, Machine, Method, Material, Measurement, Environment]
% TODO: Tambahkan Gambar Fishbone Diagram

\subsection{Teknik Pengumpulan Data}
\label{subsec:teknik-pengumpulan-data}
% [ISI DI SINI — Data Log + Kuesioner SUS + Observasi]

\subsection{Teknik Analisis Data}
\label{subsec:teknik-analisis-data}
% [ISI DI SINI — Statistik Deskriptif + K-Means + Fishbone + Thematic]

\subsection{Metode Pengujian}
\label{subsec:metode-pengujian}
% [ISI DI SINI — Black-Box + Akurasi Model + SUS]

% ------------------------------------------------------------
\section{DESAIN SISTEM}
\label{sec:desain-sistem}
% ------------------------------------------------------------
% TODO Phase 2: Sub-sections detail scope Dias
% Estimasi: ~12-15 halaman total
% Konsolidasi 15 section existing jadi 4 main sections dengan subsections

\subsection{Perancangan Sistem Global (IPO)}
\label{subsec:sistem-global}
% TODO Phase 2: Copy dari existing 3.1 Perancangan Sistem Global
% Estimasi: ~1.5 halaman
% Konten: Input → 4 Layer (Interactive Sim, AI Classification, Decision & Consequence, Data & Analysis) → Output
% [ISI DI SINI + Gambar 3.x Desain Sistem Global]

\subsection{Arsitektur Sistem Hybrid}
\label{subsec:arsitektur}
% TODO Phase 2: Copy dari existing 3.3 Perancangan Arsitektur Sistem
% Estimasi: ~1.5 halaman
% Konten: Browser-based inference + server-side services, Node.js/Express, SQLite
% [ISI DI SINI + Gambar 3.x Arsitektur Sistem]

\subsection{Model CNN MobileNet}
\label{subsec:cnn-model}
% TODO Phase 2: Copy dari existing 3.4 + 3.4.1-3.4.3
% Estimasi: ~2 halaman
% Konten: pemilihan kategori dataset, arsitektur, training, konversi TF.js
% [ISI DI SINI + Tabel 3.x Hyperparameter + Gambar 3.x Arsitektur CNN]

\subsection{Skema Data Log dan Backend API}
\label{subsec:data-log-api}
% TODO Phase 2: Konsolidasi dari existing 3.7 + 3.8
% Estimasi: ~2 halaman
% Konten: struktur JSON log (Tabel), endpoint REST API (Tabel), alur logging
% [ISI DI SINI + Tabel 3.x Struktur JSON Log + Tabel 3.x Endpoint REST API]

\subsection{K-Means Clustering}
\label{subsec:kmeans-clustering}
% TODO Phase 2: Copy dari existing 3.9
% Estimasi: ~2 halaman
% Konten: feature engineering, penentuan k (Elbow/Silhouette), interpretasi cluster
% [ISI DI SINI + Gambar 3.x K-Means Pipeline]

\subsection{Admin Dashboard}
\label{subsec:admin-dashboard}
% TODO Phase 2: Copy dari existing 3.10 + tambah wireframe mockup
% Estimasi: ~1.5 halaman
% Konten: fitur dashboard, alur penggunaan, mockup dashboard
% [ISI DI SINI + Gambar 3.x Wireframe Admin Dashboard]

\subsection{Pembagian Tugas dan Alur Pengerjaan}
\label{subsec:pembagian-tugas}
% TODO Phase 2: ADD section baru — DUO PATTERN (dari Daffa 3.2.2)
% Estimasi: ~1.5 halaman
% Diagram pembagian tugas (BIRU=Can, ABU=Dias)
% Narasi: penulis (Dias) ↔ rekan (Can)
% Titik temu: Probe UI
% [ISI DI SINI + Gambar 3.x Pembagian Tugas]

\subsection{Kontrak Data Dias--Can}
\label{subsec:kontrak-data}
% TODO Phase 2: Copy dari existing 3.14 + rapikan
% Estimasi: ~1 halaman
% Konten: format JSON payload, endpoint, batas tanggung jawab
% [ISI DI SINI + Gambar 3.x Kontrak Data]

\subsection{Rancangan Pengujian}
\label{subsec:rancangan-pengujian}
% TODO Phase 2: Black-Box + akurasi model (confusion matrix) + SUS untuk dashboard
% Estimasi: ~1.5 halaman
% [ISI DI SINI]

% ------------------------------------------------------------
\section{JADWAL PENELITIAN}
\label{sec:jadwal-penelitian}
% ------------------------------------------------------------
% TODO Phase 2: Tabel timeline SINKRON dengan Can
% Estimasi: ~1 halaman
% Format: 1 tabel shared, kolom Bulan 1-12, rows kegiatan + inisial C/D
% Lihat Style_Guide_v1.md Section 7.3 untuk daftar kegiatan
%
% [ISI DI SINI + Tabel 3.x Timeline]

% ============================================================
% DAFTAR PUSTAKA
% ============================================================
\chapter*{DAFTAR PUSTAKA}
\addcontentsline{toc}{chapter}{DAFTAR PUSTAKA}

% TODO Phase 2: Daftar Pustaka format PENS-IEEE
% Urut by nomor [1], [2], ... sesuai urutan muncul
% Format: Author. (Year). Title. Venue, vol. X, no. Y, pp. Z-W. doi: ...
%
% Referensi Dias (22 total — lihat proposal-dias.tex existing):
% [1] Ng et al. - AI literacy
% [2] Casal-Otero et al. - AI literacy K-12
% [3] Khosravi et al. - XAI in education
% [4] Mosqueira-Rey et al. - HITL ML state of the art
% [5] Memarian & Doleck - HITL in AIED
% [6] Xu et al. - Deep learning for free-hand sketch
% [7] Goh et al. - Front-end deep learning web apps
% [8] Smilkov et al. - TensorFlow.js
% [9] Karran et al. - Confidence visualization
% [10] Pansri et al. - K-Means clustering student learning
% [11] Alzahrani et al. - K-Means student engagement
% [12] Videnovik et al. - GBL CS education
% [13] Chan, Wan & King - Flow Theory
% [14] Schroeder et al. - Pedagogical agents
% [15] Meng et al. - MediaPipe (NOTE: bisa DEPRECATED bila Can tidak pakai finger tracking login)
% [16] Pressman - Software Engineering
% [17] Sommerville - Software Engineering
% [18] Sugiyono - R&D
% [19] Creswell - Mixed Methods
% [20] Ishikawa - Quality Control
% [21] Brooke - SUS
% [22] Bangor et al. - SUS empirical
% [ISI DI SINI]

\end{document}
